\chapter*{Remerciements}
\thispagestyle{empty}

\setlength{\parindent}{0pt}
\setlength{\parskip}{0.9em}

Je tiens tout d’abord à exprimer ma profonde gratitude à mon encadrant universitaire, Monsieur Imed Zghal, pour la qualité de son accompagnement, sa disponibilité constante et la richesse de ses conseils. Son suivi attentif et son expertise ont largement contribué à la réussite de ce projet de fin d’études.

J’adresse également mes sincères remerciements aux membres du jury pour avoir accepté d’évaluer mon travail et pour l’intérêt qu’ils lui ont accordé.

Ma reconnaissance s’étend à la Société Tunisienne de l’Électricité et du Gaz (STEG), qui m’a offert l’opportunité de réaliser mon projet de fin d’études dans un environnement professionnel aussi stimulant qu’enrichissant.

Je tiens à remercier tout particulièrement Monsieur Sabeur Abdelkafi, mon encadrant au sein de la société, pour son accompagnement, ses conseils judicieux et son soutien constant tout au long de ma période de stage. Son expérience et ses orientations m’ont permis d’acquérir de nouvelles compétences et de consolider mes connaissances.

Mes vifs remerciements vont également à l’ensemble des enseignants de l’École Supérieure de Commerce de Sfax pour la qualité de leur enseignement et la rigueur scientifique qu’ils m’ont transmise durant mon parcours universitaire.

Enfin, je tiens à remercier toutes les personnes qui, de près ou de loin, ont contribué à la réalisation de ce travail. Leur soutien, leurs encouragements et leur bienveillance ont été essentiels dans l’accomplissement de ce projet.
 

\clearpage